
\documentclass{article}


% \documentclass[tikz, margin=3mm]{standalone}


% \usepackage{draftwatermark}
% \SetWatermarkText{Draft}
% \SetWatermarkScale{5}
% \SetWatermarkLightness {0.9} 
% \SetWatermarkColor[rgb]{0.7,0,0}

\usepackage{tikz}
\usepackage{geometry}                		% See geometry.pdf to learn the layout options. There are lots.
\geometry{letterpaper}                   		% ... or a4paper or a5paper or ... 
%\geometry{landscape}                		% Activate for for rotated page geometry
%\usepackage[parfill]{parskip}    		% Activate to begin paragraphs with an empty line rather than an indent
\usepackage{graphicx}				% Use pdf, png, jpg, or eps� with pdflatex; use eps in DVI mode
								% TeX will automatically convert eps --> pdf in pdflat						
								% TeX will automatically convert eps --> pdf in pdflatex		
\usepackage{amssymb}
\usepackage{mathrsfs}
\usepackage{hyperref}
\usepackage{url}
\usepackage{subcaption}
\usepackage{authblk}
\usepackage{amsmath}
\usepackage{mathtools}
\usepackage{graphicx}
\usepackage[export]{adjustbox}
\usepackage{fixltx2e}
\usepackage{hyperref}
\usepackage{alltt}
\usepackage{color}
\usepackage[utf8]{inputenc}
\usepackage[english]{babel}
\usepackage{float}
\usepackage{bigints}
\usepackage{braket}
\usepackage{siunitx}
\usetikzlibrary{angles, quotes}
\usepackage{relsize}


%
% so you can do e.g., \begin{bmatrix}[r] (or [c] or [l])
%

\makeatletter
\renewcommand*\env@matrix[1][c]{\hskip -\arraycolsep
  \let\@ifnextchar\new@ifnextchar
  \array{*\c@MaxMatrixCols #1}}
\makeatother

\newcommand{\argmax}{\operatornamewithlimits{argmax}}
\newcommand{\argmin}{\operatornamewithlimits{argmin}}

\title{A Few Notes on the Principle of Least Action \\ and Newton's Second Law of Motion}
\author{David Meyer \\ dmm@1-4-5.net}

\date{Last update: \today}							% Activate to display a given date or no date


\begin{document}
\maketitle

\section{Introduction}
The principle of least action, which is also known as the stationary-action principle, is a variational principle that, when applied to the action of a mechanical system, 
amazingly yields the equations of motion for that system. The principle states that the trajectories (i.e. the solutions of the equations of motion) are stationary points of the 
system's action functional  \cite{malham2016}. One of the many amazing aspects of the principle is that it can be used to derive Newton's Second Law, Lagrangian and 
Hamiltonian equations of motion, and even  general relativity (see Einstein–Hilbert action). These notes explore the connection between the Principle of Least Action and 
Newton's Second Law.

\section{Using Calculus to Find The Minimum of a Function}
\label{sec:minimum}
The story starts with how we use calculus to find the minimum of a function. In particular, suppose we have a function $f(x)$. We know that if we move a 
small distance from $x$, say $\epsilon$, we know that $x \to \epsilon \Rightarrow f(x) \to f(x + \epsilon)$. We also know that we can evaluate $f(x + \epsilon)$ 
using a Taylor series \cite{wiki:taylor}. In particular


\begin{equation*}
\begin{array}{lllll}
f(x + \epsilon)
&=&  f(x) + \epsilon \dfrac{df}{dx} + \epsilon^2 \dfrac{d^2 f}{dx^2} +  \cdots                               &\qquad  \qquad \mathrel{\#} \text{definition of the Taylor series for $f(x + \epsilon)$}
\end{array}
\end{equation*}

\bigskip
\noindent
We also know that 

\begin{equation*}
\begin{array}{lllll}
\Delta f
&=& f(x + \epsilon) - f(x)                                                                                                                & \qquad \qquad \mathrel{\#} \text{definition of $\Delta$}  \\
[10pt]
&=&  f(x) + \epsilon \dfrac{df}{dx} + \epsilon^2 \dfrac{d^2 f}{dx^2} +  \cdots  - f(x)                       &\qquad  \qquad \mathrel{\#} \text{substitute Taylor series for $f(x + \epsilon)$} \\
[10pt]
&=& \epsilon \dfrac{df}{dx} + \epsilon^2 \dfrac{d^2 f}{dx^2} +  \cdots                                           &\qquad  \qquad \mathrel{\#} \text{$f(x)$ cancels} \\
[10pt]
&=& 0 + \epsilon^2 \dfrac{d^2 f}{dx^2} +  \cdots                                                                           &\qquad  \qquad \mathrel{\#} \text{$\dfrac{df}{dx} = 0$ at $f$'s minimum} \\
[10pt]
&=&  \epsilon^2 \dfrac{d^2 f}{dx^2}  + \cdots                                                                                &\qquad  \qquad \mathrel{\#} \text{simplify} \\
[10pt]
&\approx&  \epsilon^2 \dfrac{d^2 f}{dx^2}                                                                                     &\qquad  \qquad \mathrel{\#} \text{$\epsilon$ is close to zero ignore so we can higher order terms}
\end{array}
\end{equation*}

\bigskip
\noindent
So we can say that $\Delta f \sim \epsilon^2$, that is, the difference is a second order effect. In summary we know that for normal functions $f$ at the minimum

\begin{equation*}
\Delta f =  f(x + \epsilon) - f(x)  = 0
\label{eqn:Delta}
\end{equation*}

\noindent
to first order in $\epsilon$.

\bigskip
\noindent
An important (or maybe, the most important) difference between the approach taken by the Principle of Least Action and Newton's Second Law is that the Principle 
of Least Action looks at the value of an entire  \emph{path}, where as Newton's Second Law looks at points along the path. As we will see, these
two approaches turn out to be (amazingly) equivalent.

\section{The Principle of Least Action}
\label{sec:pola}
The action, denoted $\mathcal{S}$,  is defined to be the difference between a particle's kinetic energy \cite{wiki:kinetic_energy} and its potential energy \cite{wiki:potential_energy}, summed at every 
point along its path.  The Principle of Least  Action states that the path for which the sum of these differences has the minimum sum is the actual path the particle will take. To study this problem
we rely on the calculus of variation.

\bigskip
\noindent
Consider the setup shown in Figure \ref{fig:pola} and imagine that the actual path is represented by the function $x(t)$. That is, $x(t)$ represents
the path of a particle that starts out at position $x_1$ at time $t_1$ and ends up at $x_2$ at time $t_2$. Then what we want to do is to look at 
another path that is some small deviation away from $x(t)$ (we know from Equation \ref{eqn:Delta} that this will be close to $x(t)$). We hope that if 
we look at the change in the action of the path with the small deviation that change should be close to zero as well.

\bigskip
\noindent
As we can see from Figure \ref{fig:pola} we call the small deviation $\eta(t)$ and we call the deviated path $x(t) + \eta(t)$ (the red dashed curve). 

% Now recall
% that the action $\mathcal{S}$ is defined at the time integral from $t_1$ to $t_2$, that is
%
% \bigskip
% \begin{equation}
% \mathcal{S} = \int_{t_1}^{t_2} (KE - PE) \; dt
% \label{eqn:S}
% \end{equation}



\begin{figure}[H]
\centering
  \resizebox{0.65  \textwidth}{!} {                                                                                                                                                 % resize figure if you want
    \begin{tikzpicture}  [every edge quotes/.append style = {anchor=south, sloped}]
       \draw[thick, ->] (-0.35,0) -- (13,0) coordinate [label={[xshift=0.30cm, below]  $t$}] (t);                                                        % xshift -- get some extra space
       \draw[thick, ->] (0,-0.35) -- (0,11) coordinate [label={[xshift=-0.30cm, left]  $x$}] (x);                                                          % xshift -- get some extra space
       \coordinate[label=below left:$ O $] (O); 
%
% draw legend
%        
       \draw (7,8) node[right,rectangle] {$x(t)$}; 
       \draw [very thick] (9.5,8) -- (11,8);
       \draw (7,7)  node [right,rectangle] {${x(t) + \eta(t)}$};
       \draw [very thick,red,dashed] (9.5,7) -- (11,7);

%
% make some coordinates for later use
%
       \coordinate (x2) at (0,5);
       \coordinate (x1) at (0,2);
       \coordinate (t2) at (6,0);
       \coordinate (t1) at (2,0);
       \coordinate (t1x1) at (2,2);
       \coordinate (t2x2) at (6,5);
         
       \draw[draw=none] (x2) node[left] {$x_2$};
       \draw [] (t2) node[below] {$t_2$};
       \draw[draw=none] (x1) node[left] {$x_1$};     
       \draw[draw=none] (t1) node[below] {$t_1$};    
%
% draw points
%
         \fill (x2) circle(0.05);      
         \fill (x1) circle(0.05);                                                                                                                                                   % draw dots on axes
         \fill (t2) circle(0.05); 
         \fill (t1) circle(0.05);  
         \fill (t1x1) circle(0.05);   
         \fill (t2x2) circle(0.05);       
         
%
% the coordinates of the parabolas are unfortunately pretty much eyeballed
%
         
        \draw[yshift=-2cm, very thick,black] (2, 4) parabola[bend={+(0, 3)}, bend pos=0.5] (6, 7);                                       % draw parabolas
        \draw[yshift=-2cm, dashed, very thick,red] (2, 4) parabola[bend={+(0, 5)}, bend pos=0.5] (6, 7);     
        \coordinate (p1) at (4,6.5);                                                                                                                                         % coordinates on curves 
        \coordinate (p2) at (4,8.5);
        \draw[dashed,thick, <->] (p1) -- (p2) coordinate [label={[midway, right]  $\eta(t)$}];                                                 % draw \eta(t) between curves                                                                              
      
 %
 % draw the S integral in a box to the right
 %        
       \draw (9.5,2.5) node[draw,thick,rectangle] {$\mathlarger {\mathlarger {{\displaystyle \mathcal{S} = \int_{t_1}^{t_2} (KE - PE) \; dt}}}$}; 
    \end{tikzpicture}
  }                                                                                                                                                                                             % end resizebox                                                                                           
 \caption{Principle of Least Action Setup}
 \label{fig:pola}
\end{figure}

\bigskip
\noindent
Now, consider a particle moving in one dimension a constant gravitational field $g$ with potential energy $V(x)$ (also less formally called $PE$) In this case 
$V(x) = m \cdot g \cdot x$ (note that here $x$ is shorthand for $x(t)$). Then

\bigskip
\begin{equation*}
\begin{array}{lllll}
\mathcal{S}
&=& {\displaystyle \mathlarger {\int_{t_1}^{t_2}} (KE - PE) \; dt}           & \qquad \mathrel{\#} \text{definition of $\mathcal{S}$}  \\
[14pt]
&=& {\displaystyle \mathlarger{ \int_{t_1}^{t_2}} \left [ \frac{1}{2} m \left ( \dfrac{dx}{dt} \right )^2 - V(x) \right  ] dt}   
                                                                                                              & \qquad \mathrel{\#}  KE = \frac{1}{2} m \left ( \dfrac{dx}{dt} \right )^2 = \frac{1}{2} m v^2, \; PE = V(x) = m \cdot g \cdot x \\
\end{array}
\end{equation*}

\bigskip
\bigskip
\noindent
Now define $\delta \mathcal{S} = \mathcal{S}(x + \eta) - \mathcal{S}(x)$. We know that $\delta \mathcal{S} = 0$ to the first order in $\eta$ (here $\eta$ is shorthand for $\eta(t)$). So
now we can expand $\mathcal{S}(x + \eta)$ as follows


\bigskip
\begin{equation*}
\begin{array}{lllll}
\mathcal{S}(x + \eta)
&=& {\displaystyle \mathlarger{ \int_{t_1}^{t_2}} \left [ \frac{1}{2} m \left ( \dfrac{d}{dt} (x+\eta) \right )^2 - V(x+\eta) \right  ] dt}   & \qquad \mathrel{\#} \text{substitute $x + \eta$ for $x$}  \\
[12pt]
&=& {\displaystyle \mathlarger{ \int_{t_1}^{t_2}} \left [ \frac{1}{2} m \left ( \dfrac{dx}{dt} + \dfrac{d\eta}{dt} \right )^2 - V(x+\eta) \right  ] dt}   & \qquad \mathrel{\#} \text{derivative is a linear operator} 
\end{array}
\end{equation*}

\bigskip
\noindent
Next notice that 


\bigskip
\begin{equation*}
\begin{array}{lllll}
 \left ( \dfrac{dx}{dt} + \dfrac{d\eta}{dt} \right )^2 
&=&  \left (\dfrac{dx}{dt} \right)^2 + \left (\dfrac{d\eta }{dt} \right)^2  +2 \left ( \dfrac{dx}{dt} \right ) \left ( \dfrac{d\eta}{dt} \right )          &\qquad \mathrel{\#} \text{multiply out} \\
[12pt]
&=&  \left (\dfrac{dx}{dt} \right)^2 + 2 \left ( \dfrac{dx}{dt} \right ) \left ( \dfrac{d\eta}{dt} \right )                                                              
                                                                                                                                        &\qquad \mathrel {\#} \text{$\left ( \dfrac{d\eta }{dt} \right )^2 \approx 0$ is a second order term (Section \ref{sec:minimum})}
\end{array}
\end{equation*}

\bigskip
\noindent
Next want to expand $ V(x + \eta)$ using a Taylor series:

\medskip
\begin{equation*}
\begin{array}{lllll}
 V(x + \eta)
&=& V(x) + \eta \dfrac{dV}{dx} + \dfrac{\eta^2}{2!}  \dfrac{d^2V}{dx^2} + \cdots     &\qquad \mathrel{\#} \text{expand Taylor series} \\
\end{array}
\end{equation*}

\bigskip
\noindent
So now if we substitute in these expressions and group terms we see that

\bigskip
\begin{equation*}
\begin{array}{lllll}
\mathcal{S}(x + \eta)
&=& {\displaystyle \mathlarger{ \int_{t_1}^{t_2}} \left [ \frac{1}{2} m \left ( \dfrac{dx}{dt} \right )^2 - V(x) 
+ \dfrac{1}{2} m \left (2 \dfrac{dx}{dt} \dfrac{d \eta}{dt} \right ) - \eta \dfrac{dV}{dx} + \text{higher order terms}  \right ] dt} 
\end{array}
\end{equation*}

\bigskip
\noindent
We also have the constraint that

\begin{equation*}
\delta \mathcal{S} = \mathcal{S}(x+\eta) - \mathcal{S}(x) = 0
\end{equation*}

\bigskip
\noindent
and we saw earlier that

\begin{equation*}
\mathcal{S}(x) = {\displaystyle \mathlarger{ \int_{t_1}^{t_2}} \left [ \frac{1}{2} m \left ( \dfrac{dx}{dt} \right )^2 - V(x) \right  ] dt}   
\end{equation*}


\bigskip
\noindent
Putting the last three equations together we find that

\bigskip
\begin{equation*}
\delta \mathcal{S} =  {\displaystyle \mathlarger{ \int_{t_1}^{t_2}}  
\left [  m \left (\dfrac{dx}{dt} \dfrac{d \eta}{dt} \right ) - \eta \dfrac{dV}{dx} \right ] dt }  = 0
\end{equation*}

\medskip
\bigskip
\noindent
What we would like is to have this integral be zero regardless of the value of $\eta$. So how can we "factor out" $\eta$? The first thing to notice is that
by the product rule \cite{wiki:product_rule} we have

\bigskip
\begin{equation*}
\begin{array}{lllll}
\dfrac{d}{dx} \left ( \dfrac{dx}{dt} \eta \right ) 
&=&  \dfrac{d^2x}{dt^2} \eta + \dfrac{dx}{dt} \dfrac{d \eta}{dt}                                    & \mathrel{\#} \text{product rule: $ {\displaystyle {\dfrac {d}{dx}}(u\cdot v)={\dfrac {du}{dx}}\cdot v+u\cdot {\dfrac {dv}{dx}}}$} \\
[12pt]
&\Rightarrow&  \dfrac{dx}{dt} \dfrac{d \eta}{dt} \ = \dfrac{d}{dx} \left ( \dfrac{dx}{dt} \eta \right )  - \dfrac{d^2x}{dt^2} \eta & \mathrel{\#} \text{solve for $\dfrac{dx}{dt} \dfrac{d \eta}{dt}$} \\
[12pt]
&\Rightarrow& {\displaystyle  \int_{t_1}^{t_2} \dfrac{dx}{dt} \dfrac{d \eta}{dt} \ =  \int_{t_1}^{t_2} \dfrac{d}{dx} \left ( \dfrac{dx}{dt} \eta \right ) dt   - \int_{t_1}^{t_2}  \dfrac{d^2x}{dt^2} \eta \; dt } 
&  \mathrel{\#} \text{integrate both sides} \\
[12pt]
&\Rightarrow& {\displaystyle  \int_{t_1}^{t_2} \dfrac{dx}{dt} \dfrac{d \eta}{dt} \ =  \dfrac{dx}{dt} \eta \bigg |^{t_2}_{t_1} - \int_{t_1}^{t_2}  \dfrac{d^2x}{dt^2} \eta \; dt } & \mathrel{\#} \text{integral and derivative in $2^{nd}$ term cancel} \\
[12pt]
&\Rightarrow& {\displaystyle  \int_{t_1}^{t_2} \dfrac{dx}{dt} \dfrac{d \eta}{dt} \ =   - \int_{t_1}^{t_2}  \dfrac{d^2x}{dt^2} \eta \; dt }                   & \mathrel{\#} \text{by definition $\eta(t_1) = \eta(t_2) = 0$}
\end{array}
\end{equation*}

\bigskip
\noindent
So now we know that

\begin{equation}
 {\displaystyle  \int_{t_1}^{t_2} \dfrac{dx}{dt} \dfrac{d \eta}{dt} \ =   - \int_{t_1}^{t_2}  \dfrac{d^2x}{dt^2} \eta \; dt }  
 \label{eqn:intermediate}
\end{equation}

\bigskip
\noindent
and we saw earlier that 

\begin{equation}
\delta \mathcal{S} =  {\displaystyle \mathlarger{ \int_{t_1}^{t_2}}  
\left [  m \left (\dfrac{dx}{dt} \dfrac{d \eta}{dt} \right ) - \dfrac{dV}{dx} \eta  \right ] dt }  = 0
\label{eqn:prod}
\end{equation}

\bigskip
\noindent
so

\begin{equation*}
\begin{array}{lllll}
\delta \mathcal{S} 
&=&  {\displaystyle \mathlarger{ \int_{t_1}^{t_2}}   \left [  - m \dfrac{d^2x}{dt^2} \eta - \dfrac{dV}{dx} \eta \right ] dt }            &\qquad \mathrel{\#} \text{substitute result from Equation \ref{eqn:intermediate} into Equation \ref{eqn:prod}} \\
[12pt]
&=&  {\displaystyle \mathlarger{ \int_{t_1}^{t_2}}   \left [  - m \dfrac{d^2x}{dt^2} -  \dfrac{dV}{dx} \right ] \eta \; dt }                &\qquad \mathrel{\#} \text{factor out $\eta$} 
\end{array}
\end{equation*}

\bigskip
\noindent
Since $\delta \mathcal{S}$ must equal zero (Equation \ref{eqn:prod}) and we know $\eta(t)$ is not the constant function 0  it must be that

\bigskip
\begin{equation*}
- m \dfrac{d^2x}{dt^2} -  \dfrac{dV}{dx} = 0
\end{equation*}

\bigskip
\noindent
But now note that 

\begin{equation*}
 -  \dfrac{dV}{dx} =  m \dfrac{d^2x}{dt^2}
\end{equation*}

\bigskip
\noindent
and since ${\displaystyle -  \dfrac{dV}{dx} = F}$ and ${\displaystyle  \dfrac{d^2x}{dt^2} = a}$ we amazingly have arrived at Newton's Second Law \cite{wiki:newtons_laws}:

\bigskip
\begin{equation*}
F = ma
\end{equation*}



\section{Conclusions}
It's amazing that two such different approaches lead to the same result...

\bibliographystyle{plain}
\bibliography{/Users/dmm/papers/bib/qc}
\end{document} 

